\newpage
\begin{center}
  \textbf{\large 4. МЕТОДЫ ЗАЩИТЫ МОБИЛЬНЫХ ПРИЛОЖЕНИЙ ОТ ПЕРЕХВАТА ДАННЫХ}
\end{center}
\refstepcounter{chapter}
\addcontentsline{toc}{chapter}{4. МЕТОДЫ ЗАЩИТЫ МОБИЛЬНЫХ ПРИЛОЖЕНИЙ ОТ ПЕРЕХВАТА ДАННЫХ}

\section{Введение}
% 11 строк

В условиях стремительного роста числа кибератак, направленных на мобильные приложения, защита сетевых взаимодействий становится одной из ключевых задач для обеспечения безопасности данных пользователей. Перехват сетевого трафика, осуществляемый с использованием прокси-серверов, представляет собой распространённый метод, позволяющий злоумышленникам анализировать и потенциально эксплуатировать уязвимости в коммуникациях между клиентской и серверной частями приложений. В данной главе рассматриваются методы защиты мобильных приложений от перехвата данных, включая современные подходы к шифрованию, аутентификации и минимизации передаваемой информации. Особое внимание уделяется рекомендациям, разработанным на основе стандартов OWASP и других авторитетных источников, а также практическим мерам, позволяющим повысить устойчивость приложений к атакам типа "человек посередине". Предложенные методы направлены на устранение выявленных в предыдущих главах уязвимостей и обеспечение безопасной работы мобильных приложений в условиях современных киберугроз.

Рассмотренные в этой главе способы защиты от перехвата данных также замедлят разработку, уменьшат прозрачность системы и сделают аудит безопасности мобильных приложений сложнее. Однако, при разработке защищённых систем всё же следует следовать рекомендациям из этой главы.

\section{Рекомендации OWASP по защите мобильных приложений от перехвата данных}
% 98 строк

В условиях роста киберугроз, направленных на перехват сетевого трафика мобильных приложений, внедрение рекомендаций OWASP (Open Web Application Security Project) является ключевым шагом для повышения защищённости приложений. OWASP предоставляет комплексные стандарты и руководства, такие как OWASP Mobile Security Project и OWASP Mobile Top 10 \cite{owasp_mobile_security_project, owasp_mobile_top_10_2016}, которые помогают разработчикам и специалистам по информационной безопасности минимизировать уязвимости, связанные с перехватом данных. Основное внимание уделяется предотвращению атак типа "человек посередине", утечек данных и несанкционированного доступа. Ниже приведены рекомендации, основанные на стандартах OWASP, адаптированные для защиты мобильных приложений от перехвата сетевого трафика, с учётом анализа, проведённого в рамках данной работы.

\subsection{Внедрение сквозного шифрования}
OWASP подчёркивает необходимость использования сквозного шифрования для всех сетевых взаимодействий мобильного приложения. Это включает применение протоколов TLS (Transport Layer Security) версии 1.3 или выше для защиты данных, передаваемых между клиентом и сервером \cite{nist_sp800_52}. Шифрование должно охватывать не только конфиденциальные данные, такие как учетные данные пользователя, но и все сетевые запросы, включая метаданные. Для реализации рекомендуется использовать проверенные криптографические библиотеки, такие как OpenSSL или Bouncy Castle, избегая устаревших алгоритмов, таких как MD5 или SHA-1. Это снижает риск перехвата и расшифровки данных злоумышленниками, использующими прокси-серверы. Важно также регулярно обновлять сертификаты и следить за их сроком действия, чтобы избежать уязвимостей, связанных с истёкшими или скомпрометированными сертификатами.

\subsection{Закрепление сертификатов}
Для защиты от атак "человек посередине" OWASP рекомендует внедрять закрепление сертификатов, при котором мобильное приложение проверяет подлинность сертификата сервера, сравнивая его с заранее сохранённым или ожидаемым открытым ключом \cite{owasp_mobile_top_10_2016}. Это предотвращает использование поддельных сертификатов, даже если злоумышленник подменяет системный сертификат на устройстве. Например, в случае анализа приложения для заказа еды, описанного в работе, отсутствие проверки эмитента сертификата позволило перехватить трафик. Закрепление сертификатов усложняет такие атаки, но требует тщательной реализации, чтобы избежать проблем с обновлением сертификатов. OWASP советует использовать динамическое закрепление с возможностью обновления через безопасные каналы.

\subsection{Минимизация передаваемых данных}
OWASP акцентирует внимание на принципе минимизации данных, который подразумевает передачу только необходимой информации между клиентом и сервером \cite{owasp_mobile_security_testing_guide}. В ходе анализа приложений, таких как приложение для заказа еды, было выявлено, что избыточные данные, включая местоположение пользователя или внутренние идентификаторы, передаются без необходимости. Это увеличивает риск утечки данных при перехвате. Разработчикам следует строго определять, какие данные требуются для выполнения конкретной функции, и исключать передачу избыточной информации. Например, вместо отправки полного профиля пользователя серверу можно передавать только идентификатор сессии.

\subsection{Проверка целостности и подлинности сервера}
Для предотвращения атак MitM OWASP рекомендует реализовать строгую проверку подлинности сервера. Это включает валидацию доменного имени в сертификате и проверку цепочки доверия \cite{owasp_mobile_security_testing_guide}. В случае с приложением бизнес-клуба, описанным в работе, отсутствие проверки подлинности сервера позволило перехватывать трафик на физическом устройстве iOS. Разработчикам следует использовать системные API для проверки сертификатов и избегать отключения проверки SSL/TLS в отладочных целях, так как это может остаться в релизной версии приложения, создавая уязвимость.

\subsection{Использование безопасных протоколов аутентификации}
OWASP подчёркивает важность использования современных протоколов аутентификации, таких как OAuth 2.0 или OpenID Connect, для защиты учетных данных пользователей \cite{owasp_mobile_top_10_2016}. В ходе анализа приложения Minecraft было выявлено использование нестандартного протокола шифрования, что потребовало разработки специализированного прокси-сервера. Чтобы избежать подобных уязвимостей, OWASP рекомендует применять стандартизированные протоколы, обеспечивающие безопасную передачу токенов доступа и их регулярное обновление. Также следует избегать хранения паролей в открытом виде или передачи их в запросах без шифрования.

\subsection{Обфускация и защита клиентского кода}
OWASP советует применять обфускацию кода и другие методы защиты клиентской части приложения, чтобы затруднить обратную разработку и анализ сетевых взаимодействий \cite{owasp_mobile_security_project}. В работе подчёркивается, что обфускация кода усложняет процесс анализа, однако не является достаточной мерой защиты. Для повышения эффективности OWASP рекомендует комбинировать обфускацию с другими методами, такими как проверка целостности приложения и защита от отладки. Это особенно важно для приложений, работающих с конфиденциальными данными, таких как приложение бизнес-клуба с годовым оборотом от 100 миллионов рублей.

\subsection{Мониторинг и аудит сетевого трафика}
OWASP рекомендует внедрять системы мониторинга сетевого трафика для выявления подозрительных активностей, таких как попытки перехвата данных \cite{owasp_mobile_security_testing_guide}. В работе описано создание постоянного прокси-сервера для мониторинга трафика игры Minecraft, что позволило выявить риски утечки данных. Разработчикам следует интегрировать механизмы логирования и анализа трафика на серверной стороне, чтобы обнаруживать аномалии, такие как использование поддельных сертификатов или нестандартные запросы. Это также помогает в проведении регулярных аудитов безопасности \cite{mitre_attck_mobile}.

\subsection{Защита от фоновой передачи данных}
Анализ приложений в работе показал, что некоторые данные передаются в фоновом режиме, что может быть использовано для сбора информации без согласия пользователя. OWASP рекомендует минимизировать фоновую передачу данных и уведомлять пользователей о любых операциях, связанных с отправкой данных \cite{owasp_mobile_security_project}. Например, приложение для заказа еды передавало данные о местоположении в фоновом режиме, что создавало риск утечки. Разработчикам следует внедрять прозрачные механизмы уведомления и давать пользователям возможность отключать фоновую передачу данных.

\section{Рекомендации Центрального разведывательного управления}
% 71 строка
% 31 строка

На основе анализа документа «Network Operations Division Cryptographic Requirements v1.1» (NOD Specification 001), разработанного Центральным разведывательным управлением (ЦРУ), ниже приведены рекомендации по защите мобильных приложений от перехвата данных. Эти рекомендации ориентированы на обеспечение высокого уровня криптографической безопасности для инструментов, используемых в разведывательных операциях, и могут быть адаптированы для защиты сетевых взаимодействий мобильных приложений.

ЦРУ требует, чтобы все инструменты, использующие сетевые коммуникации, реализовывали один из криптографических наборов: Long-lived Suite, Short-lived Suite или Collection Encryption Suite, в зависимости от концепции операций. Для мобильных приложений рекомендуется адаптировать Long-lived Suite для длительного использования или Short-lived Suite для краткосрочных операций.
Для ключевого обмена рекомендуется использование алгоритмов Diffie-Hellman, Elliptic Curve Diffie-Hellman (ECDH) или RSA с минимальной длиной ключа 2048 бит (256 бит для эллиптических кривых). ECDH предпочтителен для обеспечения совершенной прямой секретности (Perfect Forward Secrecy).
Для аутентификации требуется применять TLS 1.2, Elliptic Curve DSA, DSA или RSA с ключами длиной не менее 2048 бит (256 бит для эллиптических кривых). Аутентификация через TLS 1.2 должна включать проверку сертификатов обеих сторон, исключая использование стандартных корневых центров сертификации.

Для всех видов криптографических ключей требуется двойная длина ключа в сравнении со стандартной. Так для AES, это 256 бит.
При обязательной проверке целостности также необходимо применение HMAC с использованием SHA-256, SHA-384 или SHA-512 с ключами длиной не менее 256 бит.

Для защиты данных, временно хранящихся на ненадёжных файловых системах или передаваемых через файлово-ориентированные механизмы, ЦРУ рекомендует использовать Collection Encryption Suite. Для мобильных приложений это означает шифрование данных, подпись зашифрованных данных с использованием асимметричного ключа и генерация уникального сеансового ключа для каждой передачи. Этот сеансовый ключ шифруется с испльзованием открытого ключа получателся.

ЦРУ подчёркивает важность строгого управления ключами для обеспечения уникальности ключей для каждого развёртывания приложения. Разработчикам мобильных приложений следует разработать план управления ключами, обеспечивающий уникальность и правильное использование ключей, а также их архивирование в защищённой среде.

Для предотвращения атак с повторением сообщений ЦРУ рекомендует внедрять схемы nonce, использующие монотонные номера сообщений или, в крайнем случае, временные метки UTC. Nonce не должны повторяться для одной и той же комбинации ключей шифрования и аутентификации. Для мобильных приложений это означает внедрение уникальных идентификаторов для каждого сообщения, чтобы предотвратить повторное воспроизведение перехваченных данных.

Асимметричная криптография должна использоваться только для обмена ключами и цифровых подписей, а не для шифрования больших объёмов данных, из-за высокой вычислительной нагрузки и уязвимости к криптоанализу. Для мобильных приложений это подразумевает использование асимметричных алгоритмов (например, RSA или ECDH) исключительно для установления сеансовых ключей, а симметричных алгоритмов (например, AES) — для шифрования данных.

ЦРУ рекомендует сжимать данные перед шифрованием, используя подходящие для домена алгоритмы сжатия (например, GZIP, ZLIB для общих данных). Это уменьшает объём данных, подлежащих шифрованию, и усложняет криптоанализ. Для мобильных приложений это означает внедрение сжатия данных, передаваемых по сети, с использованием системных API для обеспечения корректности.

ЦРУ подчёркивает, что шифрование транспортного уровня (например, TLS для HTTPS) должно быть дополнительным к шифрованию, определённому криптографическими наборами. Транспортное шифрование используется для маскировки трафика, но не для обеспечения конфиденциальности, так как оно уязвимо для атак типа «человек посередине». Мобильные приложения должны реализовывать внутренний криптографический слой внутри транспортного шифрования (TLS), чтобы обеспечить защиту данных даже при компрометации внешнего слоя \cite{nist_sp800_52}.

\section{Заключение главы}
% 35 строк
% 38 строк
% 56 строк
% 5 строк

Рассмотренные в главе способы защиты от перехвата сетевого трафика исчерпывающе покрывают необходимые требования при разработке программного обеспечения. Использование даже доли этих подходов позволит значительно повысить защищённость мобильного приложения.

В рекомендациях стоит уделить отдельное внимание требованиям OWASP как самому популярному ресурсу, освещающему безопасность мобильных приложений, а также пособию Центрального Разведывательного Управления, полагаясь на авторитет организации.

